\chapter{Porter et rendre traçable}
\label{ch:porter}
% BUDGET : 5 pages.

Porter un modèle n'est pas le traduire. Une traduction fidèle ligne à ligne aurait reproduit
un programme que personne, sept ans après sa dernière exécution, ne savait plus faire tourner
ni interroger. Ce que la commande demandait était différent : un modèle qu'on puisse
\emph{questionner} --- changer un plafond, ajouter un indicateur, comparer deux politiques ---
sans rouvrir le code à chaque fois. Ce chapitre décrit le programme obtenu, l'architecture qui
le rend interrogeable, le protocole qui garantit qu'il dit la même chose que l'original, et le
travail qu'il a fallu pour que sa résolution reste possible.

\section{Formulation du programme}
\label{sec:formulation}

Le modèle décide, pour chaque parcelle du territoire, quelle culture y implanter. C'est un
programme linéaire en nombres entiers (\gls{milp}) purement binaire.

Soit $P$ l'ensemble des parcelles, $C$ celui des cultures, et $E \subseteq P \times C$
l'ensemble des couples \emph{éligibles} --- ceux qu'un masque agronomique autorise
(§\ref{sec:architecture}). La variable de décision est
\[
  Y_{p,c} \in \{0,1\}, \qquad (p,c) \in E,
\]
qui vaut 1 si la parcelle $p$ porte la culture $c$. Restreindre les variables aux seuls
couples éligibles n'est pas une optimisation d'implémentation : c'est ce qui ramène
$|P| \times |C| \approx \num{2000000}$ colonnes potentielles aux $\num{\mesBinaires}$ effectivement
déclarées, et c'est la première raison pour laquelle le problème est abordable.

L'objectif est la marge brute agrégée, corrigée du risque à la manière de Markowitz :
\begin{equation}
  \max \; \sum_{(p,c) \in E} Y_{p,c} \; s_p \; m_c \; \bigl(1 - a_{f(p)} \, v_c \bigr),
  \label{eq:objectif}
\end{equation}
où $s_p$ est la surface de la parcelle (ha), $m_c$ la marge brute unitaire de la culture
(\eur/ha), $f(p)$ l'exploitation à laquelle la parcelle appartient, $a_f$ le coefficient
d'aversion au risque de cette exploitation --- déduit de son type dominant --- et $v_c$ un
coefficient de perte propre à la culture. Un avertissement s'impose ici, car le paramètre
porte dans les données un nom trompeur : $v_c$ n'est \emph{pas} une variance. C'est une
fraction de marge perdue, la pénalité est donc linéaire en surface, sans terme quadratique ni
covariance entre cultures. L'équation~\eqref{eq:objectif} reste par conséquent linéaire, et
c'est heureux --- mais elle ne réalise pas l'arbitrage moyenne-variance que son nom d'origine
suggère.

Les contraintes se rangent en cinq familles. À chaque parcelle, au plus une culture :
\begin{equation}
  \sum_{c \,:\, (p,c) \in E} Y_{p,c} \;\le\; 1, \qquad \forall p \in P.
  \label{eq:unicite}
\end{equation}
Par exploitation, des règles agronomiques de rotation et de part maximale --- l'ananas au plus
sur 75\,\% de la surface, la jachère au moins à 20\,\% de la banane export --- toutes de la
forme
$\sum_{p \in P_f} \sum_{c \in G} Y_{p,c} s_p \lesseqgtr \alpha \sum_{p \in P_f} s_p$
pour un groupe de cultures $G$ et une exploitation $f$. Par exploitation encore, un plafond de
main d'\oe uvre,
\begin{equation}
  \sum_{p \in P_f} \sum_{c} Y_{p,c} \, s_p \, h_c \;\le\; \sigma \, H_f,
  \label{eq:mo}
\end{equation}
où $h_c$ est le besoin en travail de la culture (h/ha), $H_f$ la main d'\oe uvre observée sur
l'exploitation et $\sigma$ un facteur de desserrement. Au niveau du territoire, des bornes de
production $\sum_{(p,c) \in E, c \in G} Y_{p,c} s_p y_c \lesseqgtr T$ avec $y_c$ le rendement
(t/ha) --- ce sont les quotas de filière. Enfin, des bornes portant non sur une production mais
sur un \emph{indicateur} quelconque exposé par le modèle,
\begin{equation}
  \sum_{(p,c) \in E} Y_{p,c} \, s_p \, w_p \, r_c \;\lesseqgtr\; T,
  \label{eq:indicateur}
\end{equation}
où $r_c$ est un taux par hectare (kg d'azote, \gls{ift}, heures de travail, euros de
subvention) et $w_p$ un poids de parcelle. Ce dernier détail n'est pas cosmétique : l'eau
n'est prélevée que sur les parcelles irrigables, et omettre $w_p$ fait compter
\num{56} millions de mètres cubes là où le territoire en consomme \num{35}.
L'équation~\eqref{eq:indicateur} déclinée par zone --- île, région, bassin, exploitation ---
donne la forme de la directive nitrates, avec un seuil exprimé par hectare de la zone.

Le programme résolu pour la calibration compte ainsi $\num{\mesBinaires}$ variables binaires ;
selon la politique simulée, les scénarios prospectifs restent dans la plage
\numrange{310000}{335000}.

\section{Une architecture qui sépare le générique du territoire}
\label{sec:architecture}

La formulation ci-dessus ne mentionne la Guadeloupe nulle part. Elle parle de parcelles, de
cultures et d'exploitations. Cette observation, banale une fois faite, est le principe
d'organisation du code : un noyau générique qui ne connaît ni la canne, ni la chlordécone, ni
les groupements fonciers, et une couche d'étude de cas qui apporte les données, les règles
agronomiques et les indicateurs propres au territoire. Une troisième couche, strictement en
lecture, ne fait qu'exploiter des résultats déjà écrits.

Le noyau ne définit pas des contraintes, il définit des \emph{familles} de contraintes,
enregistrées sous un nom. Poser un plafond d'azote à \SI{1351632}{\kilogram} ne demande alors
pas d'écrire du code, mais une ligne de configuration nommant la famille, l'indicateur, le
sens de l'inégalité et le seuil. Le même mécanisme sert aux objectifs, aux critères
d'éligibilité et aux règles catégorielles. La conséquence pratique est que la totalité du
dispositif expérimental du chapitre~\ref{ch:explorer} --- dix politiques, douze forçages, deux
fronts de Pareto --- est écrite en \gls{yaml} et n'a demandé aucune modification du
constructeur de modèle.

Cette séparation a un coût qu'il faut nommer. Un nom mal orthographié dans une configuration
ne devient une erreur qu'au moment où le registre est interrogé, c'est-à-dire tard ; et une
famille de contraintes dont le module n'a pas été importé est simplement absente du registre,
donc introuvable. C'est le prix ordinaire de l'indirection, et il se paie en discipline de
vérification --- ce qui est l'objet de la section~\ref{sec:gardefous}.

\section{La parité comme protocole}
\label{sec:parite}

La consigne était de reproduire le comportement du \gls{gams} d'origine, pas de l'améliorer.
Cette contrainte a paru d'abord une gêne ; elle s'est révélée être la seule façon de savoir,
à chaque désaccord entre les deux modèles, lequel des deux se trompait.

Le portage a donc été conduit équation par équation, chacune recevant un statut explicite :
\emph{portée} (un équivalent Python actif), \emph{implicite} (obtenue sans équation dédiée),
\emph{différée} (identifiée, non faite) ou \emph{écartée} (délibérément non portée). Toute la
couche d'éligibilité --- bornes numériques, règles catégorielles, vingt-quatre interdictions
géographiques d'itinéraire technique --- est portée ; quatorze équations sont implicites parce
qu'elles annulent des codes agrégats dont l'économie est nulle côté Python, si bien que le
solveur ne les choisit jamais et que l'interdiction est sans objet ; le bloc de canne fibre
est différé en entier, ses données n'étant pas lues.

Le cas intéressant est celui des \emph{bugs}. \texttt{MODELE.txt} interroge une table sur des
colonnes qu'elle ne contient pas ; \gls{gams} renvoie zéro sans avertissement, et le test
$0 \neq 1$ devient toujours vrai. La conséquence réelle est qu'une variante de verger pluvial
est interdite sur \emph{toutes} les parcelles du territoire, ce que personne n'a jamais voulu.
Un port \guillemotleft~amélioré~\guillemotright{} aurait corrigé cela en silence, et tout
écart ultérieur entre les deux modèles serait devenu indémêlable. Le choix retenu est de
porter le comportement réel, et d'écrire la variante \guillemotleft~intention
présumée~\guillemotright{} juste à côté dans la configuration, désactivée : basculer les deux
lignes suffit à tester l'autre lecture. Le mandat de parité cesse alors d'être une contrainte
subie pour devenir un protocole : il rend chaque divergence attribuable.

Un piège de portage mérite d'être signalé pour ce qu'il enseigne. Le fichier des types de sol
les énumère dans un ordre, et le \gls{gams} associe la colonne numérique correspondante dans
un ordre \emph{différent}. Indexer par position produit des coefficients faux mais
parfaitement plausibles --- des rendements crédibles sur les mauvais sols. Aucun test ne
signale une telle erreur, aucune sortie n'a l'air anormale. C'est le genre de défaut qu'on ne
trouve qu'en relisant la source, et il justifie à lui seul que le \gls{gams} d'origine soit
resté la référence tout au long du travail plutôt qu'une documentation intermédiaire.

\section{Les garde-fous}
\label{sec:gardefous}

Trois dispositifs, qui n'attrapent pas les mêmes choses.

Les tests unitaires --- $\num{\mesTests}$ fonctions réparties sur $\num{\mesTestsFichiers}$
fichiers --- exercent chaque famille de contraintes sur des configurations minuscules
construites à la main. Ils ne lisent \emph{jamais} les données réelles, et c'est délibéré :
un test qui dépend d'un jeu de données de production échoue pour des raisons qui n'ont rien à
voir avec la logique testée, et devient un test qu'on désactive. Le revers est net : ils
vérifient que chaque contrainte fait ce qu'elle dit, jamais que le résultat global est le bon.

Cette seconde question est celle du \emph{golden snapshot}. Le dispositif construit le jeu de
données réel complet et calcule tous les blocs d'indicateurs sur une allocation fixée, puis
sérialise environ $\num{\mesChecksums}$ sommes de contrôle numériques. Il ne résout aucun
\gls{milp}, il est donc exactement reproductible, et il se compare d'une exécution à l'autre :
toute dérive numérique introduite par un remaniement apparaît immédiatement, y compris là où
aucun test unitaire ne regarde.

Reste ce qu'aucun des deux ne couvre : la résolution elle-même. Un garde-fou distinct rejoue
les métriques des trois runs de référence déclarés et les compare aux valeurs attendues. Sa
tolérance est fixée à 5\,\%, et cette valeur n'est pas arbitraire : trois exécutions du même
modèle avec trois graines \gls{highs} différentes rendent des \gls{pad} territoriaux de
\num{\mesGraineA}, \num{\mesGraineB} et \num{\mesGraineC}~\%. L'arbitraire de branchement du
solveur vaut donc environ \num{\mesBruitPad} point, et une garde à l'égalité exacte échouerait
sur une simple ré-exécution.

\section{Rendre le problème traitable}
\label{sec:tractabilite}

À l'échelle du territoire complet, la résolution demande de trente à cinquante-cinq minutes.
Cette fourchette est large parce que la dispersion est réelle : à \emph{taille constante},
les vingt résolutions enregistrées s'étalent de \SI{\mesSolveMin}{\second} à
\SI{\mesSolveMax}{\second}, soit un facteur \num{6.5} et un coefficient de variation de
\SI{\mesSolveCV}{\percent}. Le goulot n'est pas l'enveloppe logicielle, négligeable, mais la
recherche par séparation et évaluation elle-même. La conséquence méthodologique est qu'aucune
estimation ponctuelle de durée n'est honnête, et l'outil de suivi affiche donc un intervalle
issu des exécutions comparables plutôt qu'un chiffre.

Plus gênant : un solve qui atteint sa limite de temps ne rend pas une solution
\emph{imprécise}, il rend une solution \emph{prouvablement fausse}. La démonstration est
directe. En partant d'une allocation existante et en basculant, vers le meilleur fruitier
éligible, les parcelles où cela coûte le moins, on construit en quelques secondes une solution
réalisable meilleure de \num{\mesIncumbentGain} millions d'euros --- soit
\SI{\mesIncumbentEcart}{\percent} --- que celle trouvée par \gls{highs} en une heure. Il ne
s'agit donc pas d'une présomption. La signature à reconnaître est caractéristique : une
contrainte portant sur l'ananas fait reculer la canne et la banane, c'est-à-dire touche des
postes qu'elle n'a aucune raison de toucher.

Le correctif est l'amorçage à chaud. Fournir au solveur une allocation réalisable issue d'une
exécution antérieure fait passer un même problème de \SI{\mesWarmAvant}{\second} à
\SI{\mesWarmApres}{\second} pour un optimum identique --- car un amorçage ne change jamais
l'optimum, seulement la vitesse à laquelle il est \emph{prouvé}. Une précaution est ici
essentielle et facile à manquer : \gls{highs} rejette silencieusement une solution initiale
non réalisable. Sans audit, une exécution paraît amorcée et se comporte exactement comme une
exécution à froid. Chaque graine est donc évaluée contre les contraintes du modèle avant le
solve, et une graine rejetée est annoncée comme telle plutôt que supposée bonne.

Une dernière source d'intraitabilité mérite mention parce qu'elle relie ce chapitre au
suivant. Vingt-cinq variantes maraîchères, censées décrire des itinéraires agroécologiques
distincts, sont \emph{identiques au bit près} sur tous les paramètres que le modèle lit :
même marge, même rendement, même besoin en travail, même azote, même empreinte d'éligibilité.
Les rouvrir ferait passer le problème de $\num{\mesBinairesRef}$ à $\num{\mesBinairesKarus}$
variables, dont plus d'un demi-million de pure symétrie --- la recherche arborescente passant
son temps à permuter des solutions équivalentes. Mais la conséquence de modélisation est plus
grave que la conséquence algorithmique : une contrainte imposant une part de production
biologique bâtie sur ces variantes est satisfaite \emph{à coût exactement nul}, en choisissant
la copie dont le nom contient \guillemotleft~bio~\guillemotright. Une politique agroécologique
qui ne contraint rien.

\begin{reflexion}
Ces vingt-cinq variantes portent le nom de la microferme expérimentale sur laquelle la lettre
de mission prévoyait que j'aille acquérir des données technico-économiques. Elles sont vides
parce que ce terrain n'a pas eu lieu --- absorbé par la reconstruction dont ce chapitre rend
compte. Le constat n'est pas une confidence ajoutée après coup : il a la forme d'un résultat
technique, mesurable dans le nombre de variables du programme et dans le fait qu'une
contrainte de bio y soit un artefact. C'est aussi ce qui m'a fait comprendre que le débogage,
dans ce travail, n'était pas une activité subalterne préalable à la modélisation, mais la
manière dont le modèle rendait compte de ce que le projet savait et de ce qu'il ignorait. Un
paramètre absent ne laisse pas de trou : il laisse un nombre plausible.
\end{reflexion}
